\section{Research Gap}
\label{sec:related:gap}

The preceding three sections taken together set out the analytical landscape that the
present project enters. The existing tools surveyed in Section
\ref{sec:related:tools} serve passenger routing; the existing research surveyed in
Sections \ref{sec:related:ops} and \ref{sec:related:viz} provides methods for parts of
the analytical problem: quality-of-service indicators from telemetry, traffic sensing
from bus GPS, short-horizon arrival prediction, visualisation primitives for
spatial-temporal data, the Multiple Coordinated Views interaction model, and the design
principles of analytical charts. No integrated analytical platform that combines all of
these methods into a single tool for the Ho Chi Minh City bus network has been published.

The gap can be decomposed into three specific open problems that the proposed system
addresses.

\begin{enumerate}
  \item \noindent\textbf{Cross-source integration.} The bus GPS feed and the AFC feed are
        separate data sources with different calendar windows and different identifier
        spaces. The existing literature treats the two sources in isolation: the
        quality-of-service work of Srisakaokul and colleagues uses GPS alone, the
        traffic-sensing work of Derevitskiy and colleagues uses GPS alone, and AFC
        ridership has been studied largely outside the operational-analytics literature.
        No published method integrates the two feeds at a granularity that supports both
        reliability questions and demand questions on the same analytical surface.
  \item \noindent\textbf{Stakeholder-aware analytical access.} The four user roles named
        in Chapter \ref{chapter:solution} (passenger, driver, operating company, and
        public-transport authority) have different analytical questions. The existing
        tools serve only the passenger role; the existing research presents results from
        a single perspective per study, typically that of a transport-research analyst.
        An analytical platform that addresses all four roles within a single
        coordinated interface has not been demonstrated for any city, let alone for
        Ho Chi Minh City.
  \item \noindent\textbf{City-specific deployment.} Most of the academic methods cited
        in the preceding sections were demonstrated on data from cities other than
        Ho Chi Minh City. The Quality of Service framework was demonstrated on Bangkok
        \cite{srisakaokul2023}; the traffic-sensing method was demonstrated on a
        European city \cite{derevitskiy2016}; the visualisation primitives were
        developed in the European context \cite{liu2022}. Applying any of these methods
        to a Ho Chi Minh City dataset requires data preprocessing, identifier
        reconciliation between the GPS and AFC feeds, geocoding of stops and points of
        interest, and parameter tuning that the published literature does not document
        for this network.
\end{enumerate}

The proposed system addresses all three open problems through the design developed in
Chapter \ref{chapter:solution}. Cross-source integration is achieved through the
layered data architecture of Section \ref{sec:solution:architecture}, which reconciles
the GPS and AFC identifier spaces and joins them at an aggregate granularity that
preserves the analytical questions of interest. Stakeholder-aware analytical access is
achieved through the analytical themes of Section \ref{sec:solution:themes} and the
interaction model of Section \ref{sec:solution:interaction}, which together cover the
analytical questions of all four user roles within a single coordinated interface.
City-specific deployment is achieved through the preprocessing pipeline of Section
\ref{sec:solution:preprocessing}, which performs the identifier reconciliation,
geocoding, and parameter tuning required by the Ho Chi Minh City dataset. The chapter
that follows develops the technical foundations on which the design rests; Chapter
\ref{chapter:solution} then develops the design itself.
